

\section{Background and Motivation}
\label{sec:intro_background}

Penetration testing -- the authorized, systematic attempt to discover and exploit weaknesses in a target system before a real adversary does -- has traditionally resisted automation because it demands the kind of open-ended, context-sensitive reasoning that scripted scanners cannot perform. Over the past two years, agentic Large Language Models (LLMs) have repeatedly closed parts of this gap. \citet{fang2024hackwebsites} showed that a single LLM agent equipped with a browser and a small set of tools could autonomously perform multi-step SQL-injection attacks against live websites without being told which vulnerability to look for. \citet{fang2024oneday} extended this to real, critical-severity CVEs, showing that GPT-4 could exploit 87\% of a 15-vulnerability benchmark when given the CVE description, though this collapsed to 7\% once the description was withheld -- the first clear evidence that \emph{exploiting} a known vulnerability and \emph{discovering} one are almost entirely different capabilities. \citet{zhu2024hptsa} pushed further into genuinely zero-day territory with a hierarchical team of planning and task-specific agents (HPTSA), and a wide subsequent literature -- surveyed systematically in Chapter~\ref{chapter_backStudy} -- has proposed increasingly elaborate multi-agent, memory-augmented, and planning-augmented architectures for autonomous vulnerability assessment and penetration testing (VAPT).

A systematic reading of this literature -- a corpus of 29 papers spanning single-agent exploitation studies, hierarchical multi-agent frameworks, classical-planning hybrids, cross-mission memory systems, and seven independent benchmark suites -- together with four independently drafted agent-architecture proposals and expert adjudication between them, is the foundation on which this report is built. That reading converges on one finding echoed independently by at least six of the surveyed systems (AWE, AutoPT, VulnBot, PentestAgent, D-CIPHER, and Incalmo): \textbf{architecture, not model scale, is the dominant variable in autonomous exploitation performance.} A well-structured pipeline running a comparatively inexpensive model consistently outperforms an unstructured ReAct-style loop running a frontier model. Yet the same reading surfaces two independent, unresolved failure modes that no single surveyed system addresses simultaneously, and it is the combination of these two gaps -- not either one alone -- that motivates the system proposed in this report.
